% Gamma Calibration Results
% Auto-generated by find_optimal_gamma.py

\subsection{Optimal Gamma Selection}

To determine the optimal covariance inflation factor (\(\gamma\)) for domain adaptation, we systematically evaluated 12 values ranging from 0.001 to 1.0.
Figure~\ref{fig:optimal_gamma} presents a comprehensive analysis across multiple criteria.

\paragraph{Experimental Setup}
We calibrated the router using 1,121 domain-specific samples, starting from warmup priors trained on 80,000 samples.
The baseline configuration (\(\gamma=1.0\), no prior weakening) achieved 53.3\% strong model usage with 0.8002 average reward.

\paragraph{Key Findings}
Our analysis reveals that \textbf{\(\gamma=0.010\)} provides the optimal balance:

\begin{itemize}
    \item \textbf{Influence Balance}: Calibration/Prior ratio of 1.401, ensuring calibration data has sufficient influence without completely discarding warmup knowledge
    \item \textbf{Policy Adaptation}: 13.6 percentage point shift in routing strategy (from 53.3\% to 67.0\% strong model usage)
    \item \textbf{Quality Preservation}: 0.8510 average reward, representing +6.4\% change from baseline
    \item \textbf{Effective Sample Size}: Reduces prior strength from 80,000 to 800 effective samples, enabling rapid adaptation
\end{itemize}

\paragraph{Comparison with Alternative Values}
Table~\ref{tab:gamma_comparison} compares the recommended \(\gamma\) against alternative values.
Values that are too large (\(\gamma \geq 0.1\)) fail to sufficiently weaken the prior, resulting in minimal adaptation.
Values that are too small (\(\gamma \leq 0.001\)) over-weaken the prior, potentially discarding valuable warmup knowledge and reducing sample efficiency.

\begin{table}[t]
\centering
\caption{Gamma Factor Comparison}
\label{tab:gamma_comparison}
\begin{tabular}{lrrrrr}
\toprule
\(\gamma\) & Eff. N & Calib/Prior & Strong \% & Reward & Conv. Rate \\
\midrule
1.000 & 80,000 & 0.014 & 53.3 & 0.8002 & 0.003009 \\
0.500 & 40,000 & 0.028 & 56.2 & 0.8020 & 0.005592 \\
0.300 & 24,000 & 0.047 & 58.8 & 0.8278 & 0.007745 \\
0.200 & 16,000 & 0.070 & 61.3 & 0.8439 & 0.009699 \\
0.150 & 12,000 & 0.093 & 63.3 & 0.8448 & 0.011315 \\
0.100 & 8,000 & 0.140 & 65.9 & 0.8448 & 0.012901 \\
\bottomrule
\end{tabular}
\end{table}

\paragraph{Practical Implications}
The optimal gamma value enables efficient domain adaptation with minimal calibration data.
By reducing the effective prior strength by 99.0\%, we allow 1,121 calibration samples to meaningfully update beliefs formed from 80,000 warmup samples.
This demonstrates the practical viability of our approach: domain adaptation requires only a small fraction (1.40\%) of the original training data.
